\documentclass[conference]{IEEEtran}
\IEEEoverridecommandlockouts
\usepackage{cite}
\usepackage{amsmath,amssymb}
\usepackage{booktabs}
\usepackage{graphicx}
\usepackage{array}
\usepackage{makecell}
\usepackage{placeins}
\usepackage{url}
\usepackage[hidelinks]{hyperref}

\newcommand{\qwen}{Qwen2.5-Math-1.5B}
\newcommand{\deepseek}{DeepSeek-R1-Distill-Qwen-1.5B}
\newcommand{\mathfive}{MATH-500}

\begin{document}

\title{Mathematical Reasoning in Small Open-Weight LLMs:\\Prompting Accuracy and Response Efficiency}

\author{\IEEEauthorblockN{Baizheng Chen}
\IEEEauthorblockA{Department of Computer Science\\
Student ID: 72510800}
}

\maketitle

\begin{abstract}
This report studies Topic 1 of the CS6493 group project: mathematical reasoning ability of large language models. We evaluate two small open-weight mathematical reasoning models, \qwen{} and \deepseek{}, on \mathfive{}, GSM8K, and AIME 2024. The method set includes four assignment-aligned baselines, Chain-of-Thought zero-shot, Chain-of-Thought few-shot, Self-Refine, and Self-Consistency, plus two extensions: Plan-and-Solve as the main prompt-only method and Tool-Integrated Reasoning as an optional tool-augmented method. TIR is also interpreted as a lightweight agentic reasoning pattern because the model can decide when to call a tool, observe the result, and continue solving. In addition to accuracy and final response length, we report total generated length and a behavior-aware joint score that gates on correctness and penalizes excessive generation, repeated answer signals, and excessive reflection.
\end{abstract}

\section{Introduction}
Large language models can produce fluent step-by-step solutions, yet mathematical reasoning remains difficult because the model must parse symbols, execute multi-step logic, and return a precise final answer. The project requirement asks us to compare prompt-based methods on \mathfive{}, GSM8K, and AIME 2024 using two small target models, with accuracy and response length as key metrics.

Our study follows this requirement but adds an efficiency perspective. A method that reaches the right answer after very long or repetitive reasoning is not equivalent to a method that reaches the same answer concisely. This distinction is especially important for Self-Consistency, Self-Refine, Plan-and-Solve, and tool-augmented methods, where one problem can trigger multiple generations or tool rounds.

The main extension is \texttt{plan\_solve}. It is a pure prompt-only method, so it is directly comparable with Chain-of-Thought (CoT), Self-Refine, and Self-Consistency. We also include \texttt{tir}, Tool-Integrated Reasoning (TIR), as a separate tool-augmented extension. TIR is reported with an explicit caveat because it changes the interaction setting by allowing Python execution, but this change is also its main research value: it moves mathematical reasoning toward an agentic loop of reasoning, tool invocation, observation, and continuation.

\section{Related Work}
CoT prompting elicits intermediate reasoning steps before the final answer \cite{wei2022cot}. Self-Consistency improves robustness by sampling multiple reasoning paths and aggregating answers \cite{wang2022selfconsistency}, while Self-Refine asks the model to critique and revise its own output \cite{madaan2024selfrefine}. Plan-and-Solve prompting first asks the model to decompose the task and then solve according to the plan, targeting missing-step errors in zero-shot reasoning \cite{wang2023plansolve}.

Tool-integrated approaches combine natural-language reasoning with executable computation. ToRA shows that tool-use trajectories can substantially improve mathematical problem solving by interleaving reasoning with external computational tools \cite{gou2023tora}. Our TIR implementation is not a trained ToRA model; it is a prompt-level tool loop that asks the model to emit Python snippets when useful. This makes TIR a small agentic workflow rather than only a longer prompt: the model alternates between deciding whether a tool is needed, producing an executable action, receiving an observation, and revising the remaining solution.

The datasets and models are also grounded in prior work. GSM8K tests grade-school mathematical word problems \cite{cobbe2021gsm8k}, while MATH provides competition-style problems \cite{hendrycks2021math}. Qwen2.5-Math is a math-specialized Qwen model family \cite{yang2024qwenmath}; DeepSeek-R1-Distill-Qwen models are dense distilled reasoning models released with DeepSeek-R1 \cite{deepseek2025r1}. Finally, recent reasoning-efficiency work motivates our reflection and answer-behavior metrics: NoWait studies the cost of explicit reflection tokens such as ``Wait'' \cite{wang2025nowait}, and Dynamic Early Exit studies answer-aware truncation for long reasoning sequences \cite{yang2025earlyexit}. We use these ideas only as metric-design motivation, not as a claim that our system implements decoding-time early exit.

\section{Methodology}
\subsection{Models and Data}
We evaluate \qwen{} and \deepseek{}. The datasets are fixed sampled subsets stored in the project: 50 samples each for \mathfive{} and GSM8K, and 30 samples for AIME 2024. The total per model-method run is therefore 130 questions.

\subsection{Prompting Methods}
Table~\ref{tab:methods} defines the six reported methods and separates the dialogue pattern from the prompt idea. \texttt{cot\_zero} and \texttt{cot\_few\_shot} are single-turn methods. \texttt{self\_consistency} is not an iterative dialogue, but it is multi-sample because it launches five independent single-turn generations. \texttt{self\_refine}, \texttt{plan\_solve}, and \texttt{tir} are multi-pass methods; TIR is also an agentic tool loop because Python execution can be inserted between model calls and the tool output becomes an observation for the next reasoning step.

\begin{table}[t]
\centering
\caption{Prompting method definitions.}
\label{tab:methods}
\resizebox{\columnwidth}{!}{
\begin{tabular}{llp{0.54\columnwidth}}
\toprule
Method & Dialogue pattern & Definition \\
\midrule
\texttt{cot\_zero} & Single-turn & Zero-shot Chain-of-Thought. The model receives the problem and is asked to solve step by step, ending with a boxed final answer. \\
\texttt{cot\_few\_shot} & Single-turn & Few-shot Chain-of-Thought. The prompt adds three worked examples before the target problem to guide format and reasoning style. \\
\texttt{self\_refine} & Multi-turn & Solve-critique-refine. The model first writes a solution, then critiques it, then rewrites a final improved solution. \\
\texttt{self\_consistency} & Multi-sample & Sample-and-vote CoT. The model samples five independent single-turn solutions; final answers are extracted, majority voted, and the first response matching the voted answer is graded. \\
\texttt{plan\_solve} & Two-turn & Plan-and-Solve. The model first writes a short numbered plan, then solves the problem by following that plan. \\
\texttt{tir} & Agentic tool-loop & Tool-Integrated Reasoning. The model may emit fenced Python code, observe the returned output block, and continue solving with that evidence. \\
\bottomrule
\end{tabular}}
\end{table}

\subsection{Metrics}
Accuracy is the primary metric. A sample is correct when the parsed prediction is mathematically equal to the normalized ground truth. Response length is separated into final response length and total generated length. Final response length is the text submitted to the grader. Total generated length includes intermediate generations: plan plus solve for Plan-and-Solve, solve plus critique plus refine for Self-Refine, all samples for Self-Consistency, and the full tool transcript for TIR. When complete token-count metadata is available, total generated tokens use that recorded generation length, including generated code and recorded reasoning text; otherwise token length is approximated as $\lceil\max(\text{characters}/4,\text{words}/0.75)\rceil$. Character length is the raw string length.

For each sample, correctness gates the final score. The efficiency term is
\begin{equation}
\text{efficiency}_i =
\text{answer}^{0.5}_i
\text{length}^{0.3}_i
\text{reflection}^{0.2}_i.
\end{equation}
The final sample score is accuracy-first:
\begin{equation}
\text{score}_i = c_i \cdot \left(0.6 + 0.4 \cdot \text{efficiency}_i\right),
\end{equation}
where $c_i \in \{0,1\}$ is the correctness indicator. Thus incorrect answers still receive zero, while correct answers retain most of their accuracy credit and use efficiency as a secondary modifier. The three factors are defined as follows.

\textbf{Length factor.} Let $L_i$ be total generated tokens. With $L_{\text{ref}}=1024$ and $\tau_L=512$, 
\begin{equation}
\text{length}_i=\exp\left(-\frac{\max(0,L_i-L_{\text{ref}})}{\tau_L}\right).
\end{equation}
This penalizes only overlong generations.

\textbf{Answer factor.} We count answer signals such as \texttt{\textbackslash boxed\{}, ``final answer is'', ``the answer is'', and ``answer:''. Let $A_i$ be this count. Let $R_i$ be the tail ratio, defined as the character ratio from the first answer signal to the end of the final response; if no answer signal exists, $R_i=0$. The answer factor is
\begin{align}
\text{count}_i &= \exp\left(-\max(0,A_i-1)\right),\\
\text{tail}_i &= 
\begin{cases}
1, & A_i < 2,\\
\exp\left(-\frac{\max(0,R_i-0.2)}{0.2}\right), & A_i \ge 2,
\end{cases}\\
\text{answer}_i &= \text{count}_i \cdot \text{tail}_i.
\end{align}
It penalizes repeated answer declarations and long continuation after an early answer.

\textbf{Reflection factor.} We count reflective cues such as ``wait'', ``however'', ``check'', ``verify'', ``rethink'', and their Chinese equivalents. Let $F_i$ be this reflection count. With reference count 1 and $\tau_F=2$,
\begin{equation}
\text{reflection}_i=\exp\left(-\frac{\max(0,F_i-1)}{2}\right).
\end{equation}
The reported joint score is the mean of $\text{score}_i$. Mean efficiency on correct samples is reported separately as a diagnostic value.

\section{Experiments}
When multiple evaluated rows are available for the same model, method, and split, the tables report the row with the highest joint score. Overall rows are selected independently from the overall summaries, rather than recomputed from the per-dataset rows.

\begin{table*}[t]
\centering
\caption{Overall results across 130 questions.}
\label{tab:overall}
\resizebox{\textwidth}{!}{
\begin{tabular}{llcccc}
\toprule
Model & Method & Accuracy & Final Tok. & Total Tok. & Joint \\
\midrule
\deepseek{} & \texttt{cot\_zero} & 0.785 & 4100.1 & 4100.1 & 0.5112 \\
\deepseek{} & \texttt{cot\_few\_shot} & 0.785 & 3534.8 & 3534.8 & 0.4996 \\
\deepseek{} & \texttt{self\_refine} & 0.792 & 362.3 & 4709.8 & 0.5766 \\
\deepseek{} & \texttt{self\_consistency} & 0.854 & 4217.5 & 21449.8 & 0.5214 \\
\deepseek{} & \texttt{plan\_solve} & 0.892 & 174.1 & 2252.7 & 0.6955 \\
\deepseek{} & \texttt{tir} & 0.908 & 167.5 & 1460.3 & 0.7140 \\
\midrule
\qwen{} & \texttt{cot\_zero} & 0.531 & 2248.2 & 2248.2 & 0.5176 \\
\qwen{} & \texttt{cot\_few\_shot} & 0.515 & 2533.6 & 2533.6 & 0.4604 \\
\qwen{} & \texttt{self\_refine} & 0.508 & 348.5 & 1001.1 & 0.4726 \\
\qwen{} & \texttt{self\_consistency} & 0.562 & 848.9 & 4524.0 & 0.4789 \\
\qwen{} & \texttt{plan\_solve} & 0.785 & 380.9 & 721.7 & 0.6570 \\
\qwen{} & \texttt{tir} & 0.800 & 371.9 & 537.1 & 0.6681 \\
\bottomrule
\end{tabular}}
\end{table*}

\begin{table*}[t]
\centering
\caption{Accuracy by dataset.}
\label{tab:dataset}
\resizebox{\textwidth}{!}{
\begin{tabular}{llccc}
\toprule
Model & Method & \mathfive{} & GSM8K & AIME 2024 \\
\midrule
\deepseek{} & \texttt{cot\_zero} & 0.900 & 0.900 & 0.400 \\
\deepseek{} & \texttt{cot\_few\_shot} & 0.920 & 0.920 & 0.333 \\
\deepseek{} & \texttt{self\_refine} & 0.880 & 0.920 & 0.433 \\
\deepseek{} & \texttt{self\_consistency} & 0.920 & 0.980 & 0.533 \\
\deepseek{} & \texttt{plan\_solve} & 0.940 & 0.980 & 0.667 \\
\deepseek{} & \texttt{tir} & 0.960 & 0.940 & 0.767 \\
\midrule
\qwen{} & \texttt{cot\_zero} & 0.560 & 0.760 & 0.100 \\
\qwen{} & \texttt{cot\_few\_shot} & 0.560 & 0.780 & 0.000 \\
\qwen{} & \texttt{self\_refine} & 0.560 & 0.740 & 0.033 \\
\qwen{} & \texttt{self\_consistency} & 0.580 & 0.840 & 0.067 \\
\qwen{} & \texttt{plan\_solve} & 0.860 & 0.960 & 0.367 \\
\qwen{} & \texttt{tir} & 0.900 & 0.940 & 0.400 \\
\bottomrule
\end{tabular}}
\end{table*}

\begin{table}[t]
\centering
\caption{Overall behavior factors on correct samples.}
\label{tab:factors}
\resizebox{\columnwidth}{!}{
\begin{tabular}{llccc}
\toprule
Model & Method & Len. & Ans. & Refl. \\
\midrule
DS & \texttt{cot\_zero} & 0.629 & 0.248 & 0.017 \\
DS & \texttt{cot\_few} & 0.727 & 0.083 & 0.039 \\
DS & \texttt{self\_refine} & 0.286 & 0.648 & 0.578 \\
DS & \texttt{self\_cons.} & 0.010 & 0.160 & 0.013 \\
DS & \texttt{plan\_solve} & 0.608 & 0.368 & 0.983 \\
DS & \texttt{tir} & 0.726 & 0.339 & 0.975 \\
\midrule
Qwen & \texttt{cot\_zero} & 1.000 & 0.918 & 0.959 \\
Qwen & \texttt{cot\_few} & 1.000 & 0.585 & 0.961 \\
Qwen & \texttt{self\_refine} & 0.912 & 0.804 & 0.914 \\
Qwen & \texttt{self\_cons.} & 0.515 & 0.758 & 0.942 \\
Qwen & \texttt{plan\_solve} & 0.963 & 0.368 & 0.979 \\
Qwen & \texttt{tir} & 0.972 & 0.368 & 0.955 \\
\bottomrule
\end{tabular}}
\end{table}

Table~\ref{tab:overall} shows that \texttt{plan\_solve} is the strongest prompt-only extension in accuracy for both models and has higher joint score than all four baselines. For \qwen{}, it improves overall accuracy from 0.531--0.562 for the prompt-only baselines to 0.785. For \deepseek{}, the updated \texttt{plan\_solve} run reaches 0.892 overall accuracy, with strong performance on all three splits. After counting full generated tokens, its total generation cost is no longer as small as final response length alone suggests.

TIR remains highly competitive, reaching the highest overall accuracy for \deepseek{} at 0.908 and a joint score of 0.7140. For \qwen{}, it reaches 0.800 accuracy and 0.6681 joint score. However, it is not directly comparable as a prompt-only method because it can use Python execution. The correct interpretation is that TIR is a tool-augmented and agentic extension showing how arithmetic and symbolic computation can improve accuracy when the model is allowed to externalize part of the solution process into executable actions, observe results, and continue reasoning. This also incurs additional generation cost from code, tool rounds, and recorded reasoning.

Table~\ref{tab:factors} explains why efficiency diagnostics are still useful under the accuracy-first score. Self-Consistency can have high accuracy, but its total generated length is much larger because it samples five independent outputs. TIR and Plan-and-Solve retain much shorter final answers, but their total token cost is higher once full generation length is counted. This supports the decision to report both final response length and total generation cost.

\section{Discussion}
The safest final extension is Plan-and-Solve. It is easy to justify pedagogically: mathematical problem solving often begins by identifying the relevant steps before executing calculations. It is also fair against other prompt-only methods because it does not use an external calculator or verifier. The updated results show strong accuracy gains for both models, although its total generation cost should be reported separately from its concise final answers.

TIR is more distinctive but requires careful framing. It should not be described as just another prompt-only baseline. It changes the setting by giving the model a Python interpreter, so its gains reflect both reasoning and tool use. This setting is worth emphasizing rather than minimizing: it is a compact agentic pattern where the model chooses an action, receives an observation, and then updates the final solution. For a mathematical reasoning project, this is especially valuable because many errors in small models are arithmetic or symbolic execution errors, exactly the cases where a tool loop can help. It also points to a broader direction beyond static prompting: small mathematical models can be paired with reliable tools to form more capable solver agents.

The behavior-aware metrics are useful because raw accuracy can hide inefficient reasoning. Self-Consistency improves accuracy through sampling, but its total generation cost is high. DeepSeek CoT runs also show extremely long final responses, which hurts efficiency even when the answer is correct. The answer and reflection factors further identify repeated final-answer cues and excessive self-correction language. These metrics should be read as diagnostic signals rather than replacements for accuracy.

\section{Conclusion}
We evaluated two small open-weight mathematical reasoning models on \mathfive{}, GSM8K, and AIME 2024. The main recommendation is to present Plan-and-Solve as the primary proposed prompt-only method because it is fair, simple, and well aligned with mathematical problem solving. TIR should be presented as an optional tool-augmented and agentic extension with strong empirical results and an explicit caveat about Python execution. The final report should emphasize both accuracy and response efficiency, separating final answer length from total generated cost.

\FloatBarrier
\bibliographystyle{IEEEtran}
\bibliography{references}

\end{document}
